% Generated draft: requires mathematical and source review.
\documentclass[11pt]{article}
\usepackage[margin=25mm]{geometry}
\usepackage{fontspec}
\setmainfont{Latin Modern Roman}
\usepackage{amsmath,amssymb,mathtools,unicode-math}
\setmathfont{Latin Modern Math}
\usepackage{enumitem,needspace,xurl,hyperref}
\setlength{\parindent}{0pt}
\setlength{\parskip}{5pt}
\setlength{\emergencystretch}{3em}
\newcommand{\N}{\mathbb{N}}
\newcommand{\Z}{\mathbb{Z}}
\newcommand{\Q}{\mathbb{Q}}
\newcommand{\R}{\mathbb{R}}
\newcommand{\C}{\mathbb{C}}
\newcommand{\F}{\mathbb{F}}
\newcommand{\A}{\mathbb{A}}
\newcommand{\PP}{\mathbb{P}}
\newcommand{\E}{\mathbb{E}}
\newcommand{\eps}{\varepsilon}
\newcommand{\dd}{\,\mathrm{d}}
\begin{document}
\section*{UnsolvedMath ID: 30000728}
\textbf{Source problem code:} catalogue entry
\subsection{Definitions and mathematical statement}
Let $D\subset\R^d$ be bounded and open, and let $m$ be an admissible positive volume. A set is quasi-open if for every $\varepsilon>0$ it becomes open after adjoining an open set of Sobolev capacity less than $\varepsilon$. For such $\Omega\subset D$, use $H^1_0(\Omega)$ in the capacitary sense: Sobolev functions vanish quasi-everywhere outside $\Omega$, meaning except on a capacity-zero set. Define the $k$-th Dirichlet eigenvalue by
\[
\lambda_k(\Omega)=\inf_{\substack{L\subset H^1_0(\Omega)\\\dim L=k}}
\ \sup_{0\ne u\in L}\frac{\int|\nabla u|^2}{\int|u|^2},
\]
with multiplicities counted. Let $\widehat\Omega$ attain
\[
\min\{\lambda_k(\Omega):\Omega\subset D\text{ quasi-open},\ |\Omega|=m\}.
\]
The source poses three independent questions in this setup:
\begin{enumerate}
\item For $k=1$, improve the bound on the size of the singular part of $\partial\widehat\Omega$. The source gives neither the existing bound nor a definition of singularity or size, and does not distinguish free boundary inside $D$ from contact with $\partial D$.
\item For $k=2$, $d=2$ and rectangular $D$, must the optimizer be symmetric? The exact reflection or other symmetry, and whether every optimizer or existence of a symmetric optimizer is intended, are not specified.
\item For $k\geq3$ and connected $D$, are the associated Dirichlet eigenfunctions Lipschitz continuous? Such eigenfunctions satisfy the weak equation $-\Delta u=\lambda u$ with $u\in H^1_0(\widehat\Omega)$. The extract does not specify the eigenvalue indices meant by ``associated'', nor local versus global Lipschitz regularity, zero extension or boundary behavior.
\end{enumerate}
These questions retain the general bounded-open container class; connectedness in the third question does not supply smoothness of $\partial D$. Since quasi-open sets are specified up to capacitary changes, a canonical representative is needed before their ordinary topological boundaries and singular subsets are discussed. No unstated boundary regularity, curvature estimate, rectangle symmetry group or eigenfunction extension theorem is added.
\subsection{Short English statement}
For volume-constrained Dirichlet-eigenvalue minimizers, sharpen first-eigenvalue boundary-singularity bounds, determine second-eigenvalue symmetry in a rectangle, and decide higher-eigenfunction Lipschitz regularity in connected containers. The three targets and their missing boundary, symmetry and regularity conventions are kept separate.
\subsection{Sources}
\begin{itemize}
\item Supplied report: \emph{Mini-Workshop: Shape Analysis for Eigenvalues} (2008), Oberwolfach Reports, DOI 10.4171/owr/2007/18, pp. 995--1026; proposed year 2007. \url{https://doi.org/10.4171/owr/2007/18}.
\item Supplied reference labelled ``Lipschitz regularity progress'': \url{https://doi.org/10.1007/s00205-014-0801-6}.
\item Local UnsolvedMath registry, permanent ID 30000728, catalogue entry; created 25 August 2026, updated 16 September 2026. Its corrected\_verified note says the setup and all three questions are explicit. The supplied 21 August 2026 literature triage labels the combined record partially solved, reports regularity results in important special cases, and says the listed boundary, symmetry and higher-eigenfunction questions are not all resolved. This dated assessment is retained without assigning a new status to each subquestion.
\end{itemize}
\end{document}
